\section{The source questions and precise qualifications}

Section~6 asks whether equality of a graph code with the orthogonal dual
of its complementary graph code forces LCD, and whether the converse
holds. The forward direction holds over every field: that equality gives
$A(J-I-A)=0$, and the common projection theorem
\cite[Theorem~1.2]{U} gives LCD. Its conclusion also identifies the exact
one-dimensional defect when both degrees vanish in the field.

The converse fails already for the four-cycle over $\mathbb F_3$.
In cyclic coordinate order its code is spanned by
$(1,0,1,0)$ and $(0,1,0,1)$. Their Gram matrix is $2I_2$, so this code is
LCD. Its complementary graph is two disjoint edges, whose adjacency
matrix is invertible. The complementary code is the whole space, whose
dual is zero, and hence cannot equal the two-dimensional cycle code.
The examples of lengths $120$ and $176$ still give the affirmative
existence answer sought in the source's introduction: they go beyond
triangular graphs and have exact complementary dual codes. These
arguments make no claim about the historical priority of either implication.

The uniqueness problem for strongly regular parameters $(253,140,87,65)$,
mentioned in Section~6.3, remains outside our result. Our $253$-coordinate
code is constructed from the specified Steiner system, and no uniqueness
of every graph with those parameters is used. Remark~3's discussion of
locating the other two $15$-dimensional irreducible modules is also not a
construction of those modules here. Their abstract existence and
classification remain the explicitly retained Brauer-table input; the
specified degree-$120$ permutation module and its particular socle are
proved above.

The asymptotic Gilbert--Varshamov result of Sendrier~\cite{Sendrier}
concerns families of LCD codes. Membership of one finite code in the LCD
class does not by itself assert a finite optimality comparison or an
asymptotic property of that single code. We retain the cited asymptotic
result with that meaning. Likewise, the comparisons to code tables in
Propositions~1--4 and Corollary~1 are historical comparisons, not updated
optimality claims. None is used to establish a dimension, distance or
full automorphism group here.

The literal and intended readings of the representation statements have
been separated throughout. The $105$ global least-degree statement is
false without an additional restriction, as witnessed by the actual
nontrivial one-dimensional quotient character. The dimension qualifier
in the $176$ main theorem is necessary. For length $253$ the printed
$231$-dimensional dual has distance six rather than eight and is reducible;
its invariant $230$-dimensional subcode already disproves irreducibility.
The main Theorem~1(d), which does not assert those erroneous dual claims,
is preserved. A matrix-algebra certificate proves irreducibility of a
specified representation; global uniqueness and minimum-degree
classifications require the separately retained table inputs.

\paragraph{Computations and formal scope.}
The finite premises are the complete specified constructions, literal
matrices and group generators, full filtration and intertwiner systems,
complete weight certificates, and support-refinement upper bounds.
Their executable files and mathematical data are inventoried in full.
Discovery with graph-isomorphism software or MeatAxe is not a replacement
for these final checks. Exact arithmetic in the declared Python, C and
GAP implementations is trusted in the finite certificates. Operational
drivers and optional weakening tests are listed separately. The nine
previously checked generic Lean declarations for projection and averaging
retain their precise restricted scope; this application is not an
end-to-end Lean formalization. No new source, question, historical novelty,
or compression credit is assigned by this author draft.
